Teams that run an application on Kubernetes must decide how to cut it into services, and with that how a fixed CPU
budget is shared among processes. Comparisons on equal hardware mostly find a monolith faster, but rarely show by
how much, why, or whether a better allocation closes the gap. We compare four deployments of one application that
differ only in topology: a monolith, four services, eight services, and eight services with an Istio sidecar mesh.
They are built as one binary whose modules are placed at deployment time, and their responses are checked to be
identical. A W3C Server-Timing header splits every request's latency into application, database, serialisation,
network, framework and edge time without clock synchronisation, and cgroup counters record each container's CPU use
and throttling. In \the\numexpr\xo{cells}+\xo{prop/cells}+\xo{work/cells}+\xo{workprop/cells}+\xo{eks/cells}\relax{} fresh deployments on a local cluster and on Amazon EKS, decomposition raised
low-load p99 latency \xo{h1/min}--\xo{h1/max} times locally and \xo{eks/h1/min}--\xo{eks/h1/max} times in the cloud,
mostly as inter-service transit of a few tenths of a millisecond per call. Under an equal CPU budget, capacity fell
to \xo{peakpct/min}--\xo{peakpct/max}\% (local) and \xo{eks/peakpct/min}--\xo{eks/peakpct/max}\% (cloud) of the
monolith's, bounded by the CPU quota of the edge service. Allocating CPU in proportion to demand raised it to
\xo{prop/peakpct/min}--\xo{prop/peakpct/max}\%, and with 1\,ms of CPU in every module operation to
\xo{workprop/peakpct/min}--\xo{workprop/peakpct/max}\%.
% TODO(results): one sentence on autoscaling.
No decomposed configuration overtook the monolith. A sidecar mesh added
\xo{mesh/perhop/min}--\xo{eks/mesh/perhop/max}\,ms per call and more CPU than the application. The performance cost
of decomposition is a price per call that shrinks as services do more work but did not vanish in any setting we
measured. The testbed, data and analysis are public.
